\section{The Worm and the Darkness}

Christ's refrain names two things that do not end, the fire and the worm, and it begins with the worm: ``Where their worm dieth not.'' The words come from the last verse of Isaias, where the prophet sees those who have transgressed against the Lord, ``their worm shall not die, and their fire shall not be quenched,''\endnote{Isaias 66:24, Douay--Rheims (Challoner).} and the Lord repeats them. Where the fire comes upon the lost from without, the worm is a pain born within them.

\subsection{The remorse that does not repent}

The common teaching of the Doctors reads the worm as the remorse of conscience. The Supplement to the \emph{Summa} puts it this way: the worm ascribed to the damned ``must be understood to be not of a corporeal but of a spiritual nature: and this is the remorse of conscience, which is called a worm because it originates from the corruption of sin, and torments the soul, as a corporeal worm born of corruption torments by gnawing.''\endnote{\emph{Summa theologiae}, Supplement, q.~97, a.~2, corpus (1922).} Augustine let each reader choose whether the worm be bodily or spiritual; but either way the soul is gnawed, for ``it is necessarily understood that in a body thus tormented the soul also is tortured with a fruitless repentance.''\endnote{Augustine, \emph{City of God} XXI.9.}

The lost see what they have done, and they grieve, but their grief can no longer bring them back. The fifth chapter of Wisdom gives their words, where the wicked behold the just whom they despised:

\begin{quote}
Therefore we have erred from the way of truth, and the light of justice hath not shined unto us, and the sun of understanding hath not risen upon us. We wearied ourselves in the way of iniquity and destruction, and have walked through hard ways, but the way of the Lord we have not known. What hath pride profited us? or what advantage hath the boasting of riches brought us? All those things are passed away like a shadow\ldots{} Such things as these the sinners said in hell.\endnote{Wisdom 5:6--9, 14, Douay--Rheims (Challoner).}
\end{quote}

Saint Cyprian, the martyr-bishop of Carthage, set this very passage before Demetrianus, a pagan who blamed the Christians for the plagues of the age, and drew this conclusion from it: ``The pain of punishment will then be without the fruit of penitence; weeping will be useless, and prayer ineffectual. Too late they will believe in eternal punishment who would not believe in eternal life.''\endnote{Cyprian, \emph{To Demetrianus} 24, trans.\ Ernest Wallis, \emph{The Ante-Nicene Fathers}, vol.~5.}

Saint Gregory Nazianzen names the shame of conscience together with the loss of God. At the judgment, he says, some will be welcomed into the vision of the Trinity; ``the others among other torments, but above and before them all must endure the being outcast from God, and the shame of conscience which has no limit.''\endnote{Gregory Nazianzen, \emph{Oration} 16.9, trans.\ C. G. Browne and J. E. Swallow, \emph{Nicene and Post-Nicene Fathers}, 2nd series, vol.~7.} Saint Catherine's \emph{Dialogue} also puts the loss first and the worm second: the first pain, the loss of the vision of God, ``revives in them, then, the second, the worm of Conscience, which gnaws unceasingly, seeing that the soul is deprived of Me.''\endnote{Catherine of Siena, \emph{Dialogue}, ``Of the four principal torments of the damned.''} The soul that sees itself deprived of God knows also that its own act shut it out.

Saint Bernard, writing to Pope Eugenius on the things a man should consider, found the worm in memory.

\begin{quote}
This is `the worm that dieth not'---the memory of the past. Once it gets within, or rather is born within through sin, there it stays, and never by any means can be plucked out. It never ceases to gnaw the conscience; feeding on it as on food that never can be consumed it prolongs the life of misery. I shudder as I contemplate this biting worm, this never-dying death. I shudder at the thought of being the victim of this living death, this dying life.

This is the second death, which never kills, but is always killing.\endnote{Bernard of Clairvaux, \emph{On Consideration} V.12.25--26, trans.\ George Lewis (Oxford: Clarendon Press, 1908), pp.~161--162.}
\end{quote}

\noindent And he gave the reason why memory cannot release the soul: ``These former things have passed away, and have not passed away. They are out of hand, but not out of mind. What has been done cannot be undone; the doing was in time, but the effect of the doing will be eternal.''\endnote{Ibid., V.12.26, p.~162.}

Alphonsus, who gave a whole consideration to the remorse of the damned, names the three thoughts that will gnaw most: the trifles for which they were lost, the little that would have saved them, and the great good they have thrown away. The second is the sharpest because it is so ordinary, and he gives it in the lost soul's own words:

\begin{quote}
If I mortified myself by not looking at such an object; if I overcame human respect at such a time; if I avoided such an occasion of sin, such a companion, such a conversation, I should not now be damned. If I had gone to confession every week; if I had frequented the congregation; if I had read every day a spiritual book; if I had recommended myself to Jesus Christ and to Mary, I should not have relapsed into sin. I have so often purposed to do these things, but have either neglected my resolutions, or, after having begun to practise these exercises, I gave them up, and therefore I am lost.\endnote{Alphonsus, \emph{Preparation for Death}, Consideration XXVIII, second point, p.~343.}
\end{quote}

\subsection{The exterior darkness}

Three times in Matthew the Lord says that the lost are cast ``into the exterior darkness,'' where there is weeping and gnashing of teeth: the children of the kingdom who did not believe, the guest without a wedding garment, the servant who buried his talent.\endnote{Matthew 8:12; 22:13; 25:30, Douay--Rheims (Challoner). The parable of the talents is the Gospel of the Thirty-third Sunday in Ordinary Time in Year A of the Postconciliar Missal.} Job, whose lament the Church reads in the Office of the Dead, speaks of hell as a house of darkness: ``If I wait hell is my house, and I have made my bed in darkness.''\endnote{Job 17:13, Douay--Rheims (Challoner); in Latin, \latin{infernus domus mea est, et in tenebris stravi lectulum meum}, the seventh lesson of Matins in the Office of the Dead, \emph{Rituale Romanum} (Ratisbon: Pustet, 1872), p.~171.} And on All Souls' Day the Church prays at the offertory of the Requiem that the souls of the faithful departed may be delivered ``from the lion's mouth, that hell engulf them not, nor they fall into darkness.''\endnote{Missal, Mass of All Souls' Day, Offertory, \latin{Domine Jesu Christe, Rex gloriae}: \latin{libera eas de ore leonis, ne absorbeat eas tartarus, ne cadant in obscurum}. Latin and English in F. X. Lasance and F. A. Walsh, \emph{The New Roman Missal} (New York: Benziger, revised edition, 1945).}

Gregory the Great, commenting on Job's ``land of misery and darkness, where the shadow of death, and no order, but everlasting horror dwelleth,'' saw a fire that burns without giving comfort. ``Here the flame that burns gives light; There, as we have shewn by the words of the Psalmist, the fire that torments veils the light.'' And he found in that place a death that cannot finish dying:

\begin{quote}
Therefore there is done upon the wretches death without death, an end without ending, failing without failing; in that both death lives, and the end is ever beginning, and the failing is unable to fail.\endnote{Gregory the Great, \emph{Morals on the Book of Job} IX.66.100, Library of the Fathers (Oxford: Parker, 1844); on Job 10:22.}
\end{quote}

\noindent He adds that the fire, which gives no light for comfort, gives just enough light for torment, that the lost may see those whom they loved against God suffering beside them. The Supplement, which cites him, concludes that in hell ``the place must be so disposed for seeing as regards light and darkness, that nothing be seen clearly, and that only such things be dimly seen as are able to bring anguish to the heart. Wherefore, simply speaking, the place is dark.''\endnote{Supplement, q.~97, a.~4, corpus; Gregory's \emph{Morals} IX is cited in the first objection and in the argument \latin{sed contra}.}

\subsection{The worm now}

Remorse gnaws in this life too, and here it is a mercy, because on this side of death it can be the beginning of contrition. The saints prayed to keep it. Alphonsus, having made the lost soul speak, turns at once to his own: ``My God, I feel a great remorse for having offended you; but this pain consoles me\ldots{} because it encourages me to hope for pardon from you, who have promised to pardon all who repent. Yes, my Lord, I repent of having outraged you. I embrace this sweet pain of remorse.''\endnote{Alphonsus, \emph{Preparation for Death}, Consideration XXVIII, first point, Affections and Prayers, pp.~341--342.}

The two kinds of remorse differ in direction. Contrition looks back at sin and turns toward Christ, as Peter did when he wept bitterly and was restored. The worm that does not die is the remorse of a soul that has refused, to the end, to let its grief become a prayer. The examination of conscience each night, the confession that names sins plainly, and the resolution made again after each failure turn remorse into contrition while there is still time.
